\addtocontents{toc}{\protect\Needspace{3\baselineskip}}
\originalchapter{8.821 课程详讲：重力涌现的条件与禁戒}
\label{ch:holo-emergence-full}
本章对应 Hong Liu 第 1 讲的全部理论线索: 为什么把重力与其他相互作用分开, 涌现重力的想法, Weinberg--Witten 两个禁戒及全息对偶改变的假设. 所依据的四页公开讲义含一页 OCW 署名页. 本章补齐论证中的表示论步骤, 不复制原讲义的宣传性评价. 前置知识是 Poincaré 对称、Noether 定理和量子态的酉变换; 能动张量的通常定义见本仓库场论卷, 曲率计算见 GR 卷.

\originalsection{背景是给定的, 还是参与动力学}
QED、弱相互作用和 QCD 通常在给定的时空背景上写作用量. 例如平直背景的 Yang--Mills 场以 $F^a_{\mu\nu}$ 及物质场描述, Lorentz 对称决定态如何分类, 路径积分定义相关函数, Wilson 重整化群说明不同能标间如何变化. 一个理论具有这些形式结构, 不意味着它在强耦合下已经可解析求解. 后面使用大 $N$ 也必须保留这个区别.

GR 的度规 $g_{\mu\nu}(x)$ 则参与动力学. 测量距离的尺、光锥和传播场共享同一个动态几何. 把度规写为背景加涨落 $g_{\mu\nu}=\eta_{\mu\nu}+\kappa h_{\mu\nu}$ 时, 线性化涨落中有无质量自旋二自由度. 由此产生一个自然问题: 普通量子系统里的复合激发, 能否恰好成为这样的无质量自旋二粒子? 若可以, 几何也许只是更基本系统的低能描述.

先比较一个熟悉例子. 固体中的声子是原子多体系统的集体激发, 微观原子排列并不需要先含一套独立声子自由度. 但声子的对称性是介质及其低能近似的对称性. 若要求一个涌现粒子同时保持\textbf{原理论相同的完整 Lorentz 对称与局域守恒能动张量}, 条件远比形成普通集体模严格. 正是这些条件产生本章的禁戒.

全息对偶提出的形式是
\[
\text{AdS 中的量子引力}
\longleftrightarrow
\text{其边界给定时空上的量子场论}.
\]
这里两边不是把各自的场简单相加; 它们拟描述同一量子系统的不同变量. 体空间的径向方向并不是边界场论多出一条已存在的普通空间坐标. 这句话稍后需由能标、局域算符和体场的字典具体化, 本章还没有得到对偶的证明.

\originalsection{无质量态与螺旋度}
取未来指向无质量四动量 $k=(E,0,0,E)$, $E>0$. 粒子的螺旋度 $\sigma$ 是沿动量方向的角动量. 绕该方向转角 $\theta$, 固定螺旋度态变成一个相位:
\[
R(\theta,\widehat{\mathbf k})|k,\sigma\rangle
=\ee^{\ii\sigma\theta}|k,\sigma\rangle.
\]
为避免主动与被动旋转约定引起无用负号争议, 以下只比较旋转权的绝对值. 无质量动量的 little group 除平面旋转外还含非紧平移部分. 普通有限螺旋度粒子在物理态上让后者平凡作用; 它在规范势极化矢量上留下的纵向变化属于规范冗余. 所以不能从四个势分量推出光子有四个物理偏振.

取两个不共线的未来无质量动量 $k,k'$. 负时间号下
\[
(k+k')^2=2k\cdot k'=-2EE'(1-\cos\vartheta)<0.
\]
故总动量是时向的, 存在它的静止系. 在该系两个动量大小相等、方向相反, 可记为 $k=(E,0,0,E)$、$k'=(E,0,0,-E)$. 这里说的是\textbf{空间总动量}为零; 四维总动量仍有能量 $2E$, 并不是四矢量 $k+k'=0$.

两个态若取同一螺旋度 $j$, 绕同一实验室 $z$ 轴旋转时, 它们相对于各自的动量方向分别是正转与反转. 因而矩阵元 $\langle k',j|\mathcal O|k,j\rangle$ 的两态相位组合带旋转权 $2j$. 这个观察把粒子的自旋条件转成了算符能否提供相同旋转权的问题.

\originalsection{四矢量流和二阶张量能携带什么权}
四矢量在绕 $z$ 旋转时, $J^0,J^z$ 不变, $J^x\pm\ii J^y$ 分别取得 $\ee^{\pm\ii\theta}$. 四矢量允许的权只有 $0,\pm1$. 若 $2j>1$, 它无法匹配上一节两态要求的权, 同螺旋度非前向矩阵元就必须为零:
\[
j>\tfrac12\quad\Longrightarrow\quad
\langle k',j|J^\mu|k,j\rangle=0
\]
对全部不共线 $k,k'$ 成立. 回到任意参考系仍为零, 因为假设 $J^\mu$ 是物理态空间上的 Lorentz 协变流.

二阶张量变换是两个四矢量变换的乘积. 两个权相加, 得到 $0,\pm1,\pm2$, 最大绝对值为 2. 因而相同推理给
\[
j>1\quad\Longrightarrow\quad
\langle k',j|T^{\mu\nu}|k,j\rangle=0.
\]
这一步还只说明非前向矩阵元消失. 必须把守恒荷与连续前向极限联系起来, 才能得到实际禁戒, 不能直接略过这一环.

对 $\langle k,\sigma|k',\sigma'\rangle=\delta_{\sigma\sigma'}\delta^{(3)}(\mathbf k-\mathbf k')$ 的归一化, 平移协变性和守恒荷给
\begin{align*}
\lim_{k'\to k}\langle k,\sigma|J^\mu|k',\sigma\rangle
&=\frac{qk^\mu}{(2\pi)^3k^0},\\
\lim_{k'\to k}\langle k,\sigma|T^{\mu\nu}|k',\sigma\rangle
&=\frac{k^\mu k^\nu}{(2\pi)^3k^0}.
\end{align*}
完整归一化证明在本卷 PS1.1 中给出. 第一式的时间分量来自 $Q$ 本征值 $q$, 第二式来自动量本征值. 如果前向极限与上述非共线极限连续, $j>1/2$ 的非零荷、或 $j>1$ 的非零能量都与旋转选择规则矛盾. 这就是两个 Weinberg--Witten 结论的核心.

\begin{proposition}
在具有物理 Lorentz 协变守恒四流、通常有限螺旋度单粒子态及所需规则前向矩阵元的理论中, 无质量 $j>1/2$ 粒子不能携带该流的非零荷. 在具有相应物理协变守恒能动张量的理论中, 不能含携带非零四动量的无质量 $j>1$ 粒子.
\end{proposition}
这个结论同时作用于基本粒子和复合粒子; 证明没有用它们的内部组成. 反过来, 一旦某个假设缺失, 不能继续把结论当作禁止所有形式重力涌现的口号.

\originalsection{三个检验例和全息的空间差别}
QED 中光子是 $j=1$, 但不带电磁 $U(1)$ 的非零电荷, 所以不违反四流结论. $SU(2)$ 规范场在 $\sigma^3/2$ 子群下可以组成带荷 $B_\mu$, 但相应颜色 Noether 流含规范势, 在物理螺旋度空间中并不满足所需协变可观测流条件; 完整候选流与补偿规范变换见 PS1.2. 这说明题目的关键词是\textbf{合格流的存在}, 而不是只看粒子自旋数字.

GR 的重力能量也不能由一个普通、局域、规范不变的 Lorentz 协变 $T^{\mu\nu}_{\rm gravity}$ 在每一点描述. 度规的微分同胚冗余使重力能量的适当陈述转向边界量、渐近荷或在指定坐标下的伪张量. 它不满足普通固定背景场论中同一套假设, 因而已知的引力子没有与定理矛盾.

若试图在普通平直四维 Lorentz 不变场论内调节耦合, 使原时空中一组复合态成为严格无质量自旋二激发, 就会遇到能动张量禁戒. QCD 的有质量自旋二共振没有这个问题, 因为定理的无质量假设不成立. 不能由存在一个有质量张量态就推断它在任意调参极限能成为引力子.

全息对偶的体引力子则不生活在原边界场论同一个时空. 边界 $d$ 维平移的守恒 $T_{\mu\nu}$ 与体 $d+1$ 维度规扰动相联系, 不是宣布边界里出现了一个满足原定理全部假设的额外普通自旋二粒子. 这里绕开了“同一时空”与“同一局域物理算符”的要求; 它只说明不存在这一直接禁戒, 仍不是对偶真实性的证明.

\begin{example}
对 $j=1$, 两态矩阵元的旋转权是 2. 四矢量最大权是 1, 因此带荷无质量矢量态受四流禁戒; 二阶张量最大权为 2, 因而同一选择规则并不禁止其能动张量矩阵元. 这与光子有能量、没有电荷相容. 对 $j=2$, 要求权 4, 二阶张量也无法提供, 才出现引力涌现需要处理的障碍.
\end{example}
下一步不直接猜想体几何的全部细节, 而分别寻找两种证据: 黑洞温度和熵为何指向面积计数, 大 $N$ 矩阵图为何组织成二维曲面的拓扑. 本课程的后续工具沿这两条明确路线建立.
